\begin{ex}
Hình chữ nhật có bao nhiêu góc vuông?
\shortans{4}
\loigiai{
Hình chữ nhật có 4 góc vuông.
}
\end{ex}

\begin{ex}
Hình vuông có bao nhiêu góc vuông?
\shortans{4}
\loigiai{
Hình vuông có 4 góc vuông.
}
\end{ex}

\begin{ex}
Dụng cụ nào trong các dụng cụ học tập thường dùng để kiểm tra góc vuông? (Nếu là ê ke ghi 1, compa ghi 2).
\shortans{1}
\loigiai{
Ê-ke dùng để kiểm tra góc vuông.
}
\end{ex}

\begin{ex}
Kim giờ và kim phút của đồng hồ tạo thành một góc vuông lúc mấy giờ đúng vào buổi sáng?
\shortans{9}
\loigiai{
Lúc 3 giờ hoặc 9 giờ tạo thành góc vuông. Lúc 9 giờ hoặc 3 giờ.
}
\end{ex}

\begin{ex}
Lúc 3 giờ đúng, góc tạo bởi kim giờ và kim phút là góc gì? (Góc vuông ghi 1, góc không vuông ghi 0).
\shortans{1}
\loigiai{
Lúc 3 giờ đúng góc tạo bởi 2 kim là góc vuông.
}
\end{ex}

\begin{ex}
Lúc 6 giờ đúng, góc tạo bởi kim giờ và kim phút có phải là góc vuông không? (Phải ghi 1, không ghi 0).
\shortans{0}
\loigiai{
Lúc 6 giờ đúng góc bẹt (không vuông).
}
\end{ex}

\begin{ex}
Một hình tam giác vuông có bao nhiêu góc vuông?
\shortans{1}
\loigiai{
Hình tam giác vuông có đúng 1 góc vuông.
}
\end{ex}

\begin{ex}
Một hình tứ giác có nhiều nhất bao nhiêu góc vuông?
\shortans{4}
\loigiai{
Một hình tứ giác có thể có nhiều nhất 4 góc vuông (như hình chữ nhật).
}
\end{ex}

\begin{ex}
Một hình tam giác có nhiều nhất bao nhiêu góc vuông?
\shortans{1}
\loigiai{
Hình tam giác chỉ có nhiều nhất 1 góc vuông.
}
\end{ex}

\begin{ex}
Tờ giấy hình chữ nhật có 4 góc. Cắt đi 1 góc thì tờ giấy mới có bao nhiêu góc?
\shortans{5}
\loigiai{
Khi cắt 1 góc của tờ giấy hình chữ nhật ta được hình ngũ giác có 5 góc.
}
\end{ex}

\begin{ex}
Một góc gồm có một đỉnh và mấy cạnh?
\shortans{2}
\loigiai{
Một góc gồm 1 đỉnh và 2 cạnh xuất phát từ đỉnh đó.
}
\end{ex}

\begin{ex}
Mỗi mặt bàn học hình chữ nhật có bao nhiêu góc vuông?
\shortans{4}
\loigiai{
Có 4 góc vuông.
}
\end{ex}

\begin{ex}
Một hình chữ nhật có 4 góc vuông, 2 hình chữ nhật có tất cả bao nhiêu góc vuông?
\shortans{8}
\loigiai{
$4 \times 2 = 8$ góc vuông.
}
\end{ex}

\begin{ex}
Trong hình vuông ABCD, có bao nhiêu góc vuông đỉnh A?
\shortans{1}
\loigiai{
Tại đỉnh A có đúng 1 góc vuông của hình vuông.
}
\end{ex}

\begin{ex}
Lúc 9 giờ đúng, góc tạo bởi kim giờ và kim phút là góc gì? (Góc vuông ghi 1, không vuông ghi 0).
\shortans{1}
\loigiai{
Lúc 9 giờ đúng là góc vuông.
}
\end{ex}

\begin{ex}
Hình tam giác có 3 góc, trong đó có 1 góc vuông thì có mấy góc không vuông?
\shortans{2}
\loigiai{
$3 - 1 = 2$ góc không vuông.
}
\end{ex}

\begin{ex}
Một hình chữ nhật bị cắt mất một góc vuông ở một đỉnh thì còn lại bao nhiêu góc vuông ban đầu?
\shortans{3}
\loigiai{
$4 - 1 = 3$ góc vuông ban đầu còn lại.
}
\end{ex}

\begin{ex}
Khung cửa sổ hình chữ nhật có mấy góc vuông?
\shortans{4}
\loigiai{
Khung cửa sổ hình chữ nhật có 4 góc vuông.
}
\end{ex}

\begin{ex}
Mỗi chiếc ê-ke có bao nhiêu góc vuông?
\shortans{1}
\loigiai{
Một chiếc ê-ke có đúng 1 góc vuông.
}
\end{ex}

\begin{ex}
3 chiếc ê-ke có tất cả bao nhiêu góc vuông?
\shortans{3}
\loigiai{
$1 \times 3 = 3$ góc vuông.
}
\end{ex}

\begin{ex}
Góc tạo bởi hai cạnh liên tiếp của quyển vở ô li là góc gì? (Góc vuông ghi 1, không vuông ghi 0).
\shortans{1}
\loigiai{
Góc của quyển vở là góc vuông.
}
\end{ex}

\begin{ex}
Một hình chữ nhật có 4 đỉnh. Mỗi đỉnh tạo thành bao nhiêu góc vuông?
\shortans{1}
\loigiai{
Mỗi đỉnh tạo thành 1 góc vuông.
}
\end{ex}

\begin{ex}
Đồng hồ chỉ 12 giờ đúng. Góc tạo bởi hai kim có phải là góc vuông không? (Phải ghi 1, không ghi 0).
\shortans{0}
\loigiai{
Lúc 12 giờ hai kim trùng nhau, góc không vuông (góc 0 độ).
}
\end{ex}

\begin{ex}
Hình ngũ giác có 5 góc. Nếu có 2 góc vuông thì có bao nhiêu góc không vuông?
\shortans{3}
\loigiai{
$5 - 2 = 3$ góc không vuông.
}
\end{ex}

\begin{ex}
Một chiếc khăn quàng đỏ hình tam giác có 1 góc tù và 2 góc nhọn. Có góc vuông nào không? (Có ghi số lượng góc vuông).
\shortans{0}
\loigiai{
Chiếc khăn quàng đỏ không có góc vuông nào.
}
\end{ex}
